
\section{Introduction}
\label{sec:intro}








Recovering 3D structure from images has traditionally relied on a Structure-from-Motion (SfM) pipeline~\citep{schoenberger2016sfm,pan2024global}. These systems decompose reconstruction into feature extraction and matching, pose estimation, triangulation, and global bundle adjustment, yielding an inherently iterative process that alternates between registering new views and re-adjusting previous estimates.

More recently, feed-forward methods~\citep{dust3r,leroy2025grounding,vggt,lin2025depth,wang2025pi,keetha2026mapanything,vggt_omega} have replaced this pipeline with an end-to-end transformer that regresses 3D geometry from images in a single forward pass.
Their gains, in line with the broader trend in computer vision~\citep{dosovitskiy2021image,zhai2022scaling,radford2021clip,oquab2023dinov2,kirillov2023sam}, have largely been driven by scaling model capacity through backbone depth and width~\citep{lin2025depth}.

Yet iteration may not have disappeared at all --- it may simply have been absorbed into network depth.
\citet{jacobs2025raptor} show that for some applications the $L$ layers of a Vision Transformer (ViT) can be replaced with $K \ll L$ recurrent applications of a looped block with little loss in accuracy.
\citet{stary2025understanding} further probe the decoder of DUSt3R~\citep{dust3r} layer-by-layer and find that its pointmap predictions are themselves iteratively refined across depth, despite the layers being independently parameterized.

Together, these observations suggest that part of the benefit of depth in modern reconstruction transformers arises from implicit iterative refinement, at the cost of redundant layer-specific parameters.



We therefore make this iterative process explicit in the architecture, rather than relying on model depth to realize it implicitly. Starting from per-view DINOv2~\citep{oquab2023dinov2} features, we apply a single looped transformer block recurrently for $K$ refinement steps. By sampling $K$ from $[K_\text{min}, K_\text{max}]$ during training, a single checkpoint exposes $K$ as an inference-time compute knob without retraining. Analyzing the trained recurrence reveals that it does not converge to a fixed point. Instead, each step progressively aligns the state's direction toward its endpoint, a regime we call \emph{directional refinement}.



The resulting model, \methodname, matches or surpasses much larger feed-forward baselines across five reconstruction benchmarks spanning indoor, outdoor, object-centric, and driving scenes, at a small fraction of their parameters and comparable or lower compute. Importantly, we show that our looping formulation with shared weights significantly outperforms an otherwise identical variant with independent per-step parameters under the same training data and compute budget. We take this as evidence that iterative refinement with shared weights is a viable alternative to parameter scaling for 3D reconstruction.

We summarize our contributions as follows:
\begin{itemize}
    \item \methodname, a looping transformer for multi-view 3D reconstruction that applies a single shared block recurrently to a DINOv2-initialized state, with each step conditioned on a continuous time interval.
    \item A variable-$K$ training recipe in which the step count is sampled per batch from $[K_\text{min}, K_\text{max}]$, yielding a single checkpoint that exposes compute as an inference-time knob.
    \item State-of-the-art reconstruction quality across five challenging benchmarks, at a small fraction of the parameters and comparable or lower compute.
\end{itemize}



